\documentclass{beamer}
\usepackage{listings}
\usepackage{booktabs}
\usepackage{bookmark}
\usepackage{makecell}
\usepackage{graphicx}
\usepackage{beamerthemesplit}
\usepackage{tikz}

\usetheme{Warsaw}

\title{Proyecto Moogle!}
\author{Jorge Alejandro Echevarría Brunet}
\institute{Facultad de Matemática y Computación\\Universidad de la Habana}
\titlegraphic{\includegraphics[width=2cm]{matcom.jpg}}
\date{}

\begin{document}

\begin{frame}
	\titlepage
\end{frame}

\section{¿ Qué es Moogle! ?}

% \begin{columns}[t]
% 	\begin{column}{.5\textwidth}
% 		\tableofcontents[sections={1-2},currentsection]
% 	\end{column}
% 	\begin{column}{.5\textwidth}
% 		\tableofcontents[sections={3-4},currentsection]
% 	\end{column}
% \end{columns}

\begin{frame}{¿ Qué es Moogle! ?}
    Moogle! es una aplicación web \textbf{totalmente original} desarrollada con tecnología
	.NET Core 7.0, específicamente usando Blazor como \textbf{framework} para la interfaz
	gráfica, y en el lenguaje C\#, cuyo propósito es buscar inteligentemente un texto
	en un conjunto de documentos.
\end{frame}









    



\section{Arquitectura Básica}
\begin{frame}{Arquitectura Básica}
	Durante el desarrollo de la aplicación fueron creadas las clases:
	\begin{itemize}
		\item Reader: Identifica los documentos y normaliza sus textos correspondientes.
		\item TF-IDF: Calcula el peso de los documentos previamente identificados.
		\item Search: Normaliza la búsqueda (query) introducida y determina los documentos más relevantes.
		\item Operators: Trabaja con los operadores utilizados en la query, influyendo directamente
		      en el puntaje de los documentos con más relevancia.
		\item Snippet: Muestra una porción de texto de cada documento relevante y una posible
		      sugerencia al usuario.
		\item Initialize: Utiliza las funciones de las clases \textbf{Reader} y \textbf{TF-IDF}
		      a la hora de lanzar la aplicación.
	\end{itemize}
\end{frame}
% Objetos mas importantes de las clases que son los diccionarios Files IDF Texts
\begin{frame}{}
La estructura de datos más importante utilizada son los diccionarios.\\
Haciendo uso de estos se realizan las operaciones pertinentes para determinar los mejores
resultados de acuerdo a la consulta realizada.
\newline
\newline
Tengamos en cuenta que los principales diccionarios a utilizar durante el Preprocesamiento son:
\begin{itemize}
	\item Files
	\item IDF
\end{itemize}
Mientras que durante la Consulta:
\begin{itemize}
	\item QueryWeight
	\item Score
\end{itemize}
% 	\item QueryWeight
% 	\item Score
% \end{itemize}
\end{frame}


\section{Flujo de Datos}
\subsection{Preprocesamiento}

\begin{frame}
	\frametitle{Flujo de Datos: Preprocesamiento}
	Al iniciarse la aplicación se hace un llamado al método \textbf{Feed} de la clase
	\textbf{Initialize}, el cual es el encargado de llamar a su vez a los métodos
	correspondientes de las clases \textbf{Reader} y \textbf{TF-IDF}, permitiendo que a la
	hora de realizar la consulta se hayan calculado los pesos de los términos en cada documento,
	así tambien como la Frecuencia Inversa de Documento (IDF) de cada uno.
	\begin{center}
		\begin{tikzpicture}[->,>=stealth,thick]
			\node (init) at (0,2) {Initialize};
			\node (files) at (3,3) {Files};
			\node (feed) at (2,2) {Feed};
			\node (idf) at (3,1) {IDF};

			\draw (init) -- (feed);
			\draw (feed) -- (files);
			\draw (feed) -- (idf);
		\end{tikzpicture}
	\end{center}

\end{frame}

\begin{frame}
	La IDF de cada término i en el total de documentos
	N está dada por la fórmula:
	$$
		IDF_{i}=\lg_{10}{(\frac{N}{n_{i}})}
	$$
	Mientras que el cálculo de pesos de cada término i en el documento j está dado por la
	fórmula:
	$$
		W_{ij}=\frac{freq_{ij}}{max_{l}freq_{lj}} IDF_{i}
	$$
\end{frame}

\subsection{Búsqueda}

\begin{frame}{Búsqueda}
	Al realizar la consulta se le aplica una normalización a sus términos y se le asocia a cada
	uno su peso mediante la función \textbf{FeedQueryWeight}, perteneciente a la clase \textbf{Search}, 
	para ello se empleó la fórmula:
	$$
		W_{ij}=(\alpha +(1-\alpha )\frac{freq_{iq}}{max_{l}freq_{lq}})\lg_{10}{(\frac{N}{n_{i}})}
	$$
	Realizada esta acción se procede a buscar entre los documentos aquellos de mayor
	relevancia mediante la función \textbf{FeedScore}, también perteneciente a la clase 
	\textbf{Search}, para determinar la relevancia de cada documento se emplea la fórmula de
	Similitud de Cosenos, dada por:
	$$
		Rsim(d_{j},q)=\frac{\sum_{i = 1}^{n} W_{ij} W_{iq}}
		{\sqrt{\sum_{i = 1}^{n} (W_{ij})^2} \sqrt{\sum_{i = 1}^{n} (W_{iq})^2}}
	$$
\end{frame}

\begin{frame}
	\begin{center}
		\begin{tikzpicture}[->,>=stealth,thick]
			\node (query) at (-3,2) {Query};
			\node (qw) at (0,2) {FeedQueryWeight};
			\node (fs) at (3,2) {FeedScore};
			\node (br) at (6,2) {Best Results};

			\draw (query) -- (qw);
			\draw (qw)    -- (fs);
			\draw (fs) 	  -- (br);
		\end{tikzpicture}
	\end{center}
	En caso de introducir algún/unos de los operadores:
	\newline
	[$\sim $ ! $\wedge$ *], estos influirán en el puntaje final de los resultados encontrados.
	
\end{frame}

\end{document}

